% ===================== CHAPITRE 4 — RÉALISATION ======================
\chapter{Réalisation}

% --- macro locale : logo d'outil non flottant, légendé et numéroté ---
\newcommand{\logofig}[3]{%  #1 = fichier (dossier logos/), #2 = hauteur, #3 = légende
\par\medskip\noindent
\begin{minipage}{\linewidth}\centering
\includegraphics[height=#2]{logos/#1}\par\smallskip
\captionof{figure}{#3}
\end{minipage}\par\medskip}

\section*{Introduction}
Ce chapitre présente la \textbf{réalisation} de l'application. Il décrit d'abord
l'environnement de travail --- matériel, logiciels et technologies mobilisés ---
puis illustre, au moyen de captures d'écran, les principales interfaces
développées.

% ----------------------------------------------------------------------
\section{Environnement de travail}

\subsection{Environnement matériel}
Le développement a été réalisé sur un poste de travail standard : ordinateur
portable doté d'un processeur Intel Core~i5, de 16~Go de mémoire vive et du
système d'exploitation Windows~11.

\subsection{Outil de rédaction du rapport}
Le présent rapport a été rédigé avec \textbf{\LaTeX}, système de composition de
documents particulièrement adapté aux écrits scientifiques et techniques.
\logofig{latex.png}{2.2cm}{Logo de « \LaTeX »}

\subsection{Environnement logiciel}
Plusieurs outils et technologies ont été mobilisés pour la réalisation du projet.

\textbf{Java (JDK~21)} : langage et plateforme d'exécution sur lesquels repose
l'ensemble de l'application.
\logofig{java.png}{2.4cm}{Logo de « Java (JDK 21) »}

\textbf{Jakarta EE~10} : ensemble de spécifications d'entreprise (Servlets, JSP,
EJB, JPA) structurant l'application en couches.
\logofig{jakartaee.png}{2.4cm}{Logo de « Jakarta EE »}

\textbf{Hibernate} : implémentation de JPA assurant le mapping objet-relationnel
entre les entités Java et la base Oracle.
\logofig{hibernate.png}{2.2cm}{Logo de « Hibernate »}

\textbf{Oracle Database} : système de gestion de base de données relationnelle
hébergeant l'ensemble des données opérationnelles et les vues décisionnelles.
\logofig{oracle.png}{1.6cm}{Logo d'« Oracle Database »}

\textbf{WildFly / JBoss} : serveur d'applications Jakarta EE sur lequel
l'application est déployée sous forme d'archive WAR.
\logofig{wildfly.png}{2.2cm}{Logo de « WildFly / JBoss »}

\textbf{Apache Maven} : outil de gestion des dépendances et d'automatisation du
\emph{build} (production du WAR).
\logofig{maven.png}{2cm}{Logo d'« Apache Maven »}

\textbf{IntelliJ IDEA} : environnement de développement intégré utilisé pour le
développement, le débogage et le déploiement de l'application.
\logofig{intellij.png}{2.4cm}{Logo d'« IntelliJ IDEA »}

\textbf{Git et GitHub} : système de gestion de versions et plateforme
d'hébergement du code source.
\par\medskip\noindent
\begin{minipage}{\linewidth}\centering

\includegraphics[height=2.2cm]{logos/git.png}\hspace{1.5cm}%

\includegraphics[height=2.2cm]{logos/github.png}\par\smallskip
\captionof{figure}{Logos de « Git » et « GitHub »}
\end{minipage}\par\medskip

\textbf{ApexCharts} : bibliothèque JavaScript de visualisation utilisée pour les
graphiques du tableau de bord embarqué.
\logofig{apexcharts.png}{2cm}{Logo d'« ApexCharts »}

\textbf{Microsoft Power BI} : outil de \emph{reporting} décisionnel connecté aux
vues analytiques de la base Oracle.
\logofig{powerbi.png}{2.2cm}{Logo de « Microsoft Power BI »}

\textbf{Python} : langage utilisé pour la chaîne \textbf{ETL} et la
\textbf{segmentation} des agences, au moyen des bibliothèques \textbf{pandas}
(manipulation des données) et \textbf{scikit-learn} (apprentissage automatique).
\par\medskip\noindent
\begin{minipage}{\linewidth}\centering

\includegraphics[height=2.2cm]{logos/python.png}\hspace{1.2cm}%

\includegraphics[height=2.2cm]{logos/pandas.png}\hspace{1.2cm}%

\includegraphics[height=2.2cm]{logos/scikitlearn.png}\par\smallskip
\captionof{figure}{Logos de « Python », « pandas » et « scikit-learn »}
\end{minipage}\par\medskip

% ----------------------------------------------------------------------
\section{Interfaces de l'application}
Cette section illustre les principales interfaces réalisées, présentées dans
l'ordre des modules fonctionnels.

\subsection{Authentification}
La figure~\ref{fig:ui_login} présente la page d'authentification : l'utilisateur y
saisit son identifiant et son mot de passe pour accéder à l'application.

\begin{figure}[htbp]
  \centering
  \includegraphics[width=\textwidth]{screens/login.png}
  \caption{Interface d'authentification}
  \label{fig:ui_login}
\end{figure}

\subsection{Gestion des agences et des employés}
Les figures~\ref{fig:ui_agences} et~\ref{fig:ui_employes} illustrent les
interfaces de gestion des agences et des employés.

\begin{figure}[htbp]
  \centering
  \includegraphics[width=\textwidth]{screens/agences.png}
  \caption{Interface de gestion des agences (sélection et consultation des membres)}
  \label{fig:ui_agences}
\end{figure}

\begin{figure}[htbp]
  \centering
  \includegraphics[width=\textwidth]{screens/employes.png}
  \caption{Interface de gestion des employés (ajout, filtres et liste)}
  \label{fig:ui_employes}
\end{figure}

\subsection{Gestion des clients et des objectifs}
Les figures~\ref{fig:ui_clients} et~\ref{fig:ui_objectifs} illustrent les
interfaces de gestion des clients et de suivi des objectifs commerciaux.

\begin{figure}[htbp]
  \centering
  \includegraphics[width=\textwidth]{screens/clients.png}
  \caption{Interface de gestion des clients (portefeuille BTK)}
  \label{fig:ui_clients}
\end{figure}

\begin{figure}[htbp]
  \centering
  \includegraphics[width=\textwidth]{screens/objectifs.png}
  \caption{Suivi des objectifs commerciaux par agence et gestionnaire}
  \label{fig:ui_objectifs}
\end{figure}

\subsection{Pointage des employés}
La figure~\ref{fig:ui_pointage} présente l'interface du module de pointage : les
cartes de synthèse (présents, retards, absents, taux de présence), le pointage du
jour (boutons « Pointer l'arrivée / le départ ») et le formulaire de saisie ou de
correction manuelle.

\begin{figure}[htbp]
  \centering
  \includegraphics[width=\textwidth]{screens/pointage.png}
  \caption{Interface du module de pointage des employés}
  \label{fig:ui_pointage}
\end{figure}

% ----------------------------------------------------------------------
\section{Déploiement de l'application}
L'application est packagée par \textbf{Maven} sous forme d'une archive web
\texttt{gestion-agences.war}, puis déployée sur le serveur d'applications
\textbf{WildFly}. La persistance repose sur une \textbf{unité de persistance JPA}
de type \texttt{JTA}, qui accède à la base Oracle via la source de données
\texttt{java:/OracleDS} déclarée sur le serveur. Le fichier \texttt{persistence.xml}
(extrait~\ref{lst:persistence}) déclare les entités du domaine et la connexion à la
base \texttt{FREEPDB1}.

\begin{lstlisting}[style=pycode,language=XML,caption={Configuration de la persistance (persistence.xml)},label={lst:persistence}]
<persistence-unit name="default" transaction-type="JTA">
    <jta-data-source>java:/OracleDS</jta-data-source>
    <class>entity.Utilisateur</class>
    <class>entity.Agence</class>
    <class>entity.BClient</class>
    <class>entity.Objectif</class>
    <class>entity.Pointage</class>
    <class>entity.Gestionnaire</class>
    <!-- ... autres entites ... -->
    <properties>
        <property name="jakarta.persistence.jdbc.driver" value="oracle.jdbc.OracleDriver"/>
        <property name="jakarta.persistence.jdbc.url"
                  value="jdbc:oracle:thin:@localhost:1521/FREEPDB1"/>
        <property name="hibernate.dialect" value="org.hibernate.dialect.OracleDialect"/>
        <property name="hibernate.hbm2ddl.auto" value="none"/>
    </properties>
</persistence-unit>
\end{lstlisting}

Le paramètre \texttt{hibernate.hbm2ddl.auto = none} indique que le schéma de la base
est géré par les scripts SQL d'initialisation (annexe) et non régénéré par
Hibernate, garantissant la maîtrise de la structure de la base en production.

% ----------------------------------------------------------------------
\section{Tests et validation}
Afin de garantir la conformité de l'application aux besoins spécifiés, des
\textbf{tests fonctionnels} ont été menés au fil du développement, module par
module. Chaque fonctionnalité livrée a été éprouvée à travers des scénarios
nominaux et des cas limites (identifiants invalides, accès non autorisé, pointage
tardif), puis validée avec l'encadrement. Le tableau~\ref{tab:tests} récapitule
les principaux scénarios de test et leur résultat.

\begin{table}[htbp]
\centering
\caption{Principaux scénarios de test fonctionnel}
\label{tab:tests}
\renewcommand{\arraystretch}{1.2}
\begin{tabularx}{\textwidth}{>{\bfseries}p{2.7cm}Xp{1.9cm}}
\toprule
\textmd{\textbf{Module}} & \textbf{Scénario testé} & \textbf{Résultat}\\
\midrule
Authentification & Connexion avec des identifiants valides puis invalides. & Conforme\\
Rôles & Accès à une fonction réservée à l'\texttt{ADMIN} par un \texttt{USER}. & Conforme\\
Agences & Création, modification et suppression d'une agence. & Conforme\\
Employés & Ajout d'un employé et rattachement à une agence. & Conforme\\
Clients & Affectation d'un client à un gestionnaire. & Conforme\\
Objectifs & Consultation de la production par agence et période. & Conforme\\
Pointage & Pointage au-delà de 9~h (détermination du statut). & Conforme (\texttt{RETARD})\\
Demandes & Validation puis refus d'une demande avec journalisation. & Conforme\\
ETL & Exécution du pipeline et production du datamart. & Conforme\\
Décisionnel & Filtrage par district dans le rapport Power BI. & Conforme\\
\bottomrule
\end{tabularx}
\end{table}

L'ensemble des scénarios a produit le comportement attendu, ce qui confirme la
\textbf{conformité fonctionnelle} de la solution aux besoins exprimés au
chapitre~2.

\section*{Conclusion}
Ce chapitre a présenté l'environnement de travail, les interfaces réalisées et la
validation fonctionnelle des modules transactionnels de l'application. Le chapitre
suivant est consacré au \textbf{volet décisionnel} : la chaîne ETL, le tableau de
bord, les rapports Power BI et la segmentation des agences.
